%%%%%%%%%%%%%%%%%%%%%%%%%%%%%%%%%%%%%%%%%%%%%%%%%%%%%%%%%%%%%%%%%%%%%%%%%%%%%%%
% CONSTANTES POUR LA PAGE DE GARDE ET LA TAILLE DE POLICE DU CODE %%%%%%%%%%%%%
%%%%%%%%%%%%%%%%%%%%%%%%%%%%%%%%%%%%%%%%%%%%%%%%%%%%%%%%%%%%%%%%%%%%%%%%%%%%%%%
\newcommand{\annee}{2026}
\newcommand{\numsujet}{15}
\newcommand{\taillecode}{\small}
\newcommand{\interlignecode}{0.85} % ratio, 1.145 pour être identique au texte
%%%%%%%%%%%%%%%%%%%%%%%%%%%%%%%%%%%%%%%%%%%%%%%%%%%%%%%%%%%%%%%%%%%%%%%%%%%%%%%
% PRÉAMBULE AVEC IMPORTATION %%%%%%%%%%%%%%%%%%%%%%%%%%%%%%%%%%%%%%%%%%%%%%%%%%
%%%%%%%%%%%%%%%%%%%%%%%%%%%%%%%%%%%%%%%%%%%%%%%%%%%%%%%%%%%%%%%%%%%%%%%%%%%%%%%
\documentclass[a4paper,12pt]{article}
\usepackage{../epnsi} % fichier epnsi.sty dans le dossier parent
%%%%%%%%%%%%%%%%%%%%%%%%%%%%%%%%%%%%%%%%%%%%%%%%%%%%%%%%%%%%%%%%%%%%%%%%%%%%%%%
% CONTENU %%%%%%%%%%%%%%%%%%%%%%%%%%%%%%%%%%%%%%%%%%%%%%%%%%%%%%%%%%%%%%%%%%%%%
%%%%%%%%%%%%%%%%%%%%%%%%%%%%%%%%%%%%%%%%%%%%%%%%%%%%%%%%%%%%%%%%%%%%%%%%%%%%%%%
\begin{document}

\pagegarde

\consignes

\textbf{Cadre}

Un cabinet vétérinaire utilise une base de données pour assurer le suivi et la
gestion des clients et animaux qui y sont suivis, dans le but d’améliorer la
qualité des soins et de fournir des services utiles, que ce soit pour le
vétérinaire ou pour ses clients.

\textbf{Objectif}

La vaccination régulière des chats (une fois par an) étant importante mais
souvent oubliée par les propriétaires, le vétérinaire souhaiterait pouvoir
envoyer des messages de rappel de vaccination par SMS, deux mois avant
l’échéance (la date à laquelle la vaccination devrait avoir lieu). Il dispose
pour cela d’une plateforme d’envoi de SMS, mais il est nécessaire de préparer
les données et les messages à envoyer.

\textbf{Description de la base de données}

Le logiciel mis en œuvre par le cabinet utilise 3 tables pour enregistrer les
animaux, leur propriétaire et les consultations.
\begin{figure}[h]
    \centering
    \includegraphics[width=0.95\linewidth]{./images/diagramme_bdd.png}
    \caption*{Schéma de la base de données}
    \vspace{3mm}
    \includegraphics[width=0.95\linewidth]{./images/extrait_bdd.png}
    \caption*{Extraits de la base de données}
\end{figure}

Le système de gestion de base de données utilisé est SQLite, et il est possible
d’interagir avec la base de données en utilisant le module \verb|sqlite|,
disponible dans la bibliothèque standard de Python.

\textbf{Rappels techniques}

La connexion à la base de données se fait en précisant l’emplacement du fichier
de la base de données. Par exemple, si le fichier \verb|cabinet.sqlite| se
trouve dans le répertoire courant, on écrit :

\begin{verbatim}
DB_PATH = "cabinet.sqlite"
conn = sqlite3.connect(DB_PATH)   
\end{verbatim}

Une fois que l’on dispose d’une connexion (variable \verb|conn|), on utilise un
curseur (variable \verb|cursor|) pour exécuter des requêtes sur la base de
données et éventuellement récupérer les résultats de la requête.

\begin{verbatim}
cursor = conn.cursor()
resultats = cursor.execute(
    "SELECT id, nom FROM animal ORDER BY nom LIMIT 10;"
)
# resultats permet de récupérer les lignes de manière progressive, sous 
# forme d'une liste itérable de tuples correspondant aux champs demandés
# Pour accéder aux champs sélectionnés :
for r in resultats:
    # les valeurs renvoyées par le SELECT
    # sont accessibles par leur position
    print("id:", r[0], ", nom:", r[1])
\end{verbatim}

Lorsqu’un paramètre doit être donné dans une requête SQL, il est fourni sous
forme d’un \verb|?| dans le texte de la requête SQL, et passé en argument à la
méthode \verb|execute| sous forme d’un tuple (chaque \verb|?| correspond à un
élément du tuple, dans l'ordre). Par exemple, pour récupérer les chats du
propriétaire d’id \verb|2| :

\begin{verbatim}
resultat = cursor.execute(
    """
    SELECT nom FROM animal
    WHERE espece = 'chat' AND id_proprietaire = ?;
    """,
    (2,)
)   
\end{verbatim}

Les clauses de langage SQL à utiliser sont : \verb|SELECT|, \verb|FROM|,
\verb|JOIN|, \verb|ON|, \verb|WHERE|, \verb|ORDER BY|.

\textbf{Numéros de téléphone}

La plateforme d’envoi de SMS requiert un numéro de téléphone au format
06XXXXXXXX ou \newline 07XXXXXXXXX, 10 chiffres au total, sans espace, ni
point, ni autre caractère. Malheureusement, les téléphones ont été enregistrés
dans la base sans distinguer les téléphones mobiles des téléphones fixes, et en
utilisant différents styles comme “(0)6 12 34 56 78” ou “01.23.45.67.89”.

\textit{Rappel.} On peut utiliser la méthode \verb|isdigit| de la classe
\verb|str| pour déterminer si un caractère est un chiffre. Par exemple :

\begin{verbatim}
>>> '2'.isdigit()
True
>>> 'C'.isdigit()
False
\end{verbatim}

\begin{question}
    Écrire une fonction python \verb|normalisation_tel| qui prend un numéro de
    téléphone \verb|tel| de type \verb|str| en paramètre et qui renvoie le numéro
    de téléphone nettoyé en ne gardant que les chiffres. Exemples :

    \begin{verbatim}
>>> normalisation_tel("(0)2 12 99 90 12")
"0212999012"
>>> normalisation_tel("0.6.12.99.90.12")
"0612999012"
>>> normalisation_tel("0-6-12-9-90-1")
renvoie "06129901"   
    \end{verbatim}
    \vspace{-5mm}
    Une fonction \verb|test_normalisation_tel| est fournie.
\end{question}

\appel

\textbf{Validation des numéros de téléphone}

La fonction \verb|validation_tel(tel)| prend en paramètre un numéro de
téléphone nettoyé, et renvoie un booléen indiquant si ce numéro est un numéro
de téléphone portable valide ou non.

\begin{question}
    Écrire un jeu de tests permettant de vérifier le bon fonctionnement de la
    fonction.
\end{question}

\appel

\textbf{Détermination de la liste des chats à vacciner}

La vaccination se faisant tous les ans, le vétérinaire souhaite envoyer trois
rappels aux propriétaires. Chaque rappel sera donc envoyé aux propriétaires des
chats ayant été vaccinés pour la \textbf{dernière} fois dans l’intervalle
J -- 13 mois à J -- 10 mois.

Pour cela, on commencera par déterminer quels chats ont été vaccinés il y a
moins de 13 mois, à l’aide de la fonction
\verb|consultation_vaccination_chat(date)|.

Puis à l’aide de la fonction \verb|derniere_vaccination(consultations)|, on
déterminera la date de dernière vaccination de ces chats.

Enfin, dans la dernière étape du projet, se fera le filtrage des chats ayant été vaccinés pour la
\textbf{dernière} fois il y a plus de 10 mois.

\textbf{Interrogation de la base de données}

Pour interroger la base de données, on utilise la fonction
\verb|consultation_vaccination_-| \verb|chat(date)| qui doit récupérer toutes
les consultations de motif vaccination des chats (espece : chat) du cabinet,
pour des dates supérieures à une date, date, exprimée au format AAAAMMJJ (ce
format permet de comparer des dates en utilisant l’ordre alphabétique).

\newpage

On souhaite récupérer les champs suivants :

\begin{itemize}
    \item id de l’animal ;
    \item nom de l’animal ;
    \item telephone du propriétaire ;
    \item date de la consultation.
\end{itemize}

triés par le champ id d’animal puis date de consultation croissants.

\begin{question}
    En vous inspirant de la fonction
    \verb|proprietaires_animaux_nes_apres(date)|, écri-re la fonction
    \verb|consultation_vaccination_chat(date)|.

    Des tests sont fournis dans la fonction
    \verb|test_consultation_vaccination_chat|, vo-tre fonction doit les passer.
\end{question}

\appel

\textbf{Détermination de la date de dernière vaccination}

Pour déterminer la date de dernière vaccination, on utilise la fonction
\verb|derniere_vaccination| qui prend en paramètre une liste de consultations
avec les champs précédents (id de l’animal, nom de l’animal, id\_proprietaire de
l’animal, date de la consultation). Elle renvoie un dictionnaire associant
comme clef l’id de chaque chat avec les informations de sa \textbf{dernière}
(plus récente) consultation.

On donne par exemple ci-dessous la liste des consultations associées à Plume et
Gollum.

\begin{verbatim}
consultations = [
    (16,  'Plume', '0.6.36.96.89.83', '20241024'),
    (16,  'Plume', '0.6.36.96.89.83', '20251125'),
    (17, 'Gollum', '0.6.36.96.89.83', '20250113') 
]   
\end{verbatim}

La dernière consultation de Plume est le 20251125, et pour Gollum, c’est le
20250113 donc l'appel \verb|derniere_vaccination(consulations)| (où
\verb|consultations| est la liste précédente) \newline doit renvoyer le
dictionnaire suivant :

\begin{verbatim}
{
    16: (16,  'Plume', '0.6.36.96.89.83', '20251125'),
    17: (17, 'Gollum', '0.6.36.96.89.83', '20250113')
}  
\end{verbatim}

\begin{question}
    Cependant, la fonction ne produit pas les résultats attendus : les tests
    fournis dans la fonction \verb|test_derniere_vaccination| échouent.
    Identifier le problème rencontré, puis corriger cette fonction pour qu’elle
    passe les tests avec succès.
\end{question}

\appel

\section*{\centering Description du dossier}

Le dossier fourni au candidat sur l’ordinateur comporte les éléments suivants :

\begin{itemize}
    \item une version PDF de l’énoncé ;
    \item un code source de départ \verb|veto.py| ;
    \item une base de données \verb|cabinet.sqlite|.
\end{itemize}

\section*{Préparation de l'environnement}

Pour faire fonctionner le code fourni dans le dossier, les bibliothèques
suivantes doivent être présentes : \verb|sqlite3|.

\end{document}